\scanpage{8}
\mspage{3}
\margin{Le foncteur est pl. fid. au mieux sur une base $S$ réduite -- faux p. ex. sur les nbs duaux d'un corps (si $I\neq\varnothing$). Pour avoir un foncteur pl fidèle, il faut \unsure{se placer} dans le contexte \uline{propre}, i.e. \unsure{remplacer} $X$ par $(\widehat X,I)$ (de sorte que $X=\widehat X-I$).}
J'appelle ``courbe \inserted{relative} standart'' sur $S$, tout élément de l'image essentielle. Le système $(Y,\widetilde A,I,\sigma)$ qui lui donne naissance s'appellera le ``complété normalisé'' \inserted{(relatif)} (par abus de \deleted{système définissant} langage, puisque c'est non seulement le schéma relatif $Y$, mais $Y$ avec la structure supplémentaire $\widetilde A,I,\sigma_{\widetilde A}\ldots$).

2. \uline{Graphe associé à une courbe standart.}

Revenons au cas d'un corps de base $k$ alg. clos, pour commencer. Soit $\unsure{X}$ une courbe standart, d'où $Y,\deleted{\unsure{\widetilde A},}I,\widetilde A,\sigma_{\widetilde A}$.

Posons
\[
(7)\qquad S=\pi_0(Y)\simeq\text{ens. des comp. irréd. de }X
\]
On a alors le diagramme d'applications canoniques entre ensembles finis
\[
(8)\qquad
\diagram{\begin{tikzcd}
& I\arrow[d,"p"]\\
\widetilde A\arrow[r,"\unsure{\sigma|o}"]\arrow[d,"q"'] & S\\
A=\widetilde A/\sigma &
\end{tikzcd}}
\]
\mspage{7}
